\section{Background}
\label{sec:background}

\subsection{Code obfuscation}
Obfuscation is a family of transformations rather than a single technique, and the members of that
family differ in what they preserve. The classical taxonomy divides them into \emph{layout}
transforms, which rewrite identifiers and strip the names, comments and annotations a reader relies
on, \emph{control} transforms, which restructure the flow the computation takes, \emph{data}
transforms, which change how values are represented and computed, and \emph{preventive} transforms,
which target the analysis tooling rather than the human
reader~\cite{collberg1997taxonomy,collberg2002manufacturing}. The first three preserve observable
behaviour by construction. For every input, the transformed program returns what the original
returned. The fourth does not, and is not meant to. Anti-debugging checks, anti-virtualisation
probes, self-modifying code and environment-keyed decryption all make behaviour contingent on the
conditions under which the program runs, which is precisely their purpose. A parallel distinction is
the level at which a transform applies. Source and bytecode obfuscators leave a readable artifact,
whereas packers and virtualisation obfuscators emit a stub that reconstructs the real code at
runtime, so there is often no static representation of the program to read at
all~\cite{roundy2013packers,schrittwieser2016protecting}. Even within the semantics-preserving
subset, preservation is of input-output behaviour and not of timing, memory use or the exact
exceptions raised.

This paper works inside the semantics-preserving, source-level subset, and the restriction is
methodological rather than incidental. Because such a transform preserves behaviour by construction,
the ground truth for a behavioural question is computable from the clean parent and is independent
of any transform, so a grader can be exact and a claim of invariance is falsifiable rather than a
matter of judgement. Preventive transforms and packers admit no comparable instrument. Their
difficulty is recovering code to read, or executing safely enough to observe anything, which is a
different problem from reading code that is already in front of you. They also carry no
transform-independent gold answer, since behaviour is the thing they make conditional. Excluding
them narrows the claim this paper can make, and Section~\ref{sec:threats} states that limit.

Within the semantics-preserving subset, the three families predict what a model must recover. A
layout transform destroys evidence the model may have been relying on but leaves the computation's
shape intact. A control transform leaves the identifiers and changes the shape. A data transform
changes the arithmetic itself, through expression rewriting, string encoding or mixed
Boolean-arithmetic identities that are equivalent but opaque~\cite{zhou2007mba}. An adaptation that
generalises across all three has learned something about the computation. One that generalises within
a family has learned that family's surface.

The property that makes this subset a usable instrument is also what makes it harder than the
literature reports. The transforms \emph{compose}. An obfuscator may rename identifiers inside
flattened control flow inside rewritten expressions, and production tools do this by default rather
than as an option. Composition is not a harder version of the same test but a different one, since a
model may handle each part alone and still fail the stack. Evaluations of model comprehension under
obfuscation have nonetheless applied transforms one at a time, or treated a tool's default pipeline
as a single undifferentiated condition~\cite{nikiema2025codebarrier,rong2025obfuscation}. In neither
case is stack depth or transform composition varied as an experimental factor, so the question of
whether capability composes has not been asked. The same holds for the human-subject
evidence~\cite{nguyen2026human}.


\subsection{Predicting behaviour in two directions}
Program understanding can be probed through behaviour prediction rather than through
generation or recovery. Given a program and an input, \emph{output prediction} asks what the
program returns. Given a program and a return value, \emph{input prediction} asks for a call
that produces it. Both tasks are established, and current code models are markedly imperfect
at each even on unobfuscated code~\cite{cruxeval}. The two are not interchangeable. Output
prediction alone gives a narrow view of dynamic code reasoning and is more exposed to
contamination, which has motivated benchmarks that treat execution reasoning as inherently
dual and evaluate paired instances in both directions~\cite{dexbench}. Outside code,
training on forward and backward forms of the same problem improves performance on both,
which suggests that the paired formulation carries information that neither direction
supplies alone~\cite{revthink}.

Pairing also exposes failures that a single direction hides. Because the two directions
share a program, a gold value and a grader, an adaptation that improves one while damaging
the other cannot be credited to better reading of the program. Such asymmetry is common.
Fine-tuning on one direction of a code transformation improves the trained direction while
driving the untrained one to near zero, a failure named cognitive
specialization~\cite{cft_bidirectional}, and recovery of the damaged direction has been
attributed to the direction present in the training data rather than to the choice of
training objective~\cite{paired_attrib}. The inverse direction also differs from the forward
one in how it must be graded. Programs are many-to-one, so a return value admits many
producing calls, and an inferred input therefore has to be executed rather than matched
against a reference.

A related but distinct line of work formulates bidirectionality over the transformation
itself, pairing obfuscation with deobfuscation and treating reversibility as the test of
understanding~\cite{cft_bidirectional}. The distinction matters for what is measured.
Recovering the original source demonstrates knowledge of the transform, whereas answering
correctly about the behaviour of a transformed program demonstrates understanding of the
computation that program denotes.

\subsection{Adapting a model}
Fine-tuning a code model on a new distribution is now routinely done with low-rank
adaptation~\cite{lora}, which freezes the base weights $W$ and learns a low-rank update $\Delta W =
BA$ with $B \in \mathbb{R}^{d \times r}$, $A \in \mathbb{R}^{r \times k}$ and $r \ll d$, applied to
each target module. The practical consequence for this paper is not efficiency but
\emph{modularity}: an adapter is a small, separable object that can be trained per transform and
combined with others afterwards, and the \emph{same} architecture, rank and target modules can be
held fixed while only the training data or the combination rule changes. Every arm below is a LoRA
at the same rank on the same base model, so differences between them are differences in how the data
was used, not in capacity.

% \subsection{Baselines}
% \label{sec:baselines}
% Six arms span the ways one might spend a fixed budget of obfuscated training data. The first two do
% not adapt weights to obfuscation at all and bound the problem; the next three are the standard
% single-adapter answers; the last composes.

% \noindent\textbf{Untuned.} The instruction-tuned checkpoint, prompted directly. It is the reference
% every other arm is measured against, and on the backward task it is a surprisingly strong one.

% \noindent\textbf{In-context.} The same model with one solved example in the prompt. It tests whether
% the task format, rather than the obfuscation, is what the untuned model is failing at.

% \noindent\textbf{Clean-code LoRA.} An adapter trained on unobfuscated programs only. It is the
% control that separates \emph{learning the task} from \emph{learning the transforms}: it sees the
% prompt format and the answer distribution but never an obfuscated program, so whatever it gains is
% not transform knowledge.

% \noindent\textbf{Pooled training (breadth).} One adapter trained on the union of every training
% condition. This is the obvious way to spend the data and the one most likely to be done in practice:
% it needs no per-transform bookkeeping, no gate, and no merge step.

% \noindent\textbf{Anchored training.} Pooled training plus a consistency term. The student reads the
% obfuscated program while a frozen clean-code teacher reads its unobfuscated parent, and the loss adds
% $\lambda\,D_{\mathrm{KL}}$ between their answer-token distributions to the ordinary cross-entropy. It
% trains on exactly the rows pooled training uses and reduces to pooled training at $\lambda = 0$,
% which makes pooled training its data-matched control. The idea is to anchor the model to the one
% thing a semantics-preserving transform leaves invariant --- the behaviour of the same computation
% written plainly.

% \noindent\textbf{Mixture of specialists.} One adapter per transform, kept separate, with a learned
% gate that mixes their low-rank updates per input. Because a learned gate can be worth nothing over no
% gate at all, we also report a uniform gate and a gate frozen at random initialisation, which separate
% what the \emph{mixture} is worth from what the \emph{routing} is worth. This arm is the most
% expressive of the six and the most expensive: every expert must be resident at inference.

% \subsection{Merging specialists into one adapter}
% \label{sec:merging}
% Merging occupies the same design point as the mixture --- per-transform specialists, composed --- but
% resolves the composition \emph{before} inference, so what is deployed is a single adapter of exactly
% the base architecture, with no gate, no expert bank and no run-time cost over ordinary fine-tuning.

% The object being merged is a \emph{task vector}: for each specialist $i$, the update $\Delta_i$ its
% training produced. Naive averaging of task vectors degrades as their number grows, because updates
% that disagree cancel. TIES attributes that to two forms of interference --- redundant entries that
% carry little signal, and outright sign disagreement on a parameter --- and removes both in three
% steps: \emph{trim} each $\Delta_i$ to its largest-magnitude entries, \emph{elect} a sign per
% parameter by summing signed magnitudes across specialists, and \emph{merge} only the entries whose
% sign matches the elected one~\cite{ties}. Nothing is trained; the rule is arithmetic on weights
% already learned. DARE prepends a second sparsification of a different kind: drop a fraction $p$ of
% each $\Delta_i$'s entries at random and rescale the survivors by $1/(1-p)$, so the update is
% preserved in expectation while its redundancy is cut, on the finding that fine-tuning deltas are
% extremely redundant~\cite{dare}. We use both, as \textsc{TIES} and \textsc{DARE-TIES}, and report
% them separately --- they are not interchangeable, and the difference between them turns out to matter
% on one model in our panel.

% Two implementation choices are worth stating because they affect what is being measured. First, we
% merge \emph{in LoRA space}, combining the $A$ and $B$ factors of the specialists directly. The
% alternative offered by common tooling is to fold each adapter into the base weights and re-extract a
% low-rank adapter by SVD, which changes the rank and introduces reconstruction error before the
% merging rule is applied at all; that would confound the merge with its re-extraction. Second, the
% specialists are combined with uniform weights. A weighted merge is a tuning knob, and leaving it
% uniform keeps the comparison against pooled training --- which also weights its conditions equally by
% construction --- a comparison of \emph{composition}, not of tuning effort.